\section{ماژول توجه با فشرده‌سازی دنباله بر پایه تجزیه مقادیر منفرد (SVD Sequence Attention)}
\label{sec:module_svd_attention}

\subsection{معرفی و انگیزه}
این مکانیزم بر مبنای مبانی نظری پیشرفته در ادبیات توجه و فشرده‌سازی \cite{eckart1936approximation} توسعه یافته است.

مکانیزم توجه با فشرده‌سازی دنباله بر پایه تجزیه مقادیر منفرد \lr{(Low-Rank SVD-Compressed Sequence Attention)} با هدف غلبه بر پیچیدگی محاسباتی و حافظه‌ای درجه دو $\mathcal{O}(N^2)$ در طول دنباله‌های بلند طراحی شده است. در لایه‌های توجه کلاسیک، تشکیل ماتریس توجه تمام‌به‌تمام مستلزم ضرب داخلی میان تمامی جفت‌توکن‌های دنباله است؛ این در حالی است که فضای زیرفضای ویژگی‌های کلید و مقدار در اغلب لایه‌های ترنسفورمر دارای بعد مؤثر یا رتبه ذاتی بسیار پایینی است.

ماژول \lr{SVDLinearAttention} (ثبت‌شده با شناسه‌های \lr{svd\_attention}، \lr{svd\_seq}، \lr{svd\_sequence} و \lr{svd\_sequence\_attention}) با استفاده از تجزیه مقادیر منفرد ماتریس تصویربرداری کلید ($\mathbf{W}_K$) در زمان مقداردهی، یک زیرفضای پایدار متعامد با رتبه $r \ll N$ استخراج می‌کند. سپس با نگاشت دنباله کلید و مقدار بر روی این زیرفضای فشرده، بعد دنباله از $N$ به $r$ انقباض یافته و پیچیدگی توجه به $\mathcal{O}(N \cdot r \cdot d)$ تبدیل می‌گردد که مقیاس‌پذیری خطی نسبت به طول توالی را به ارمغان می‌آورد.

\subsection{مبانی نظری و فرمول‌بندی ریاضی}
فرض کنید $\mathbf{X} \in \mathbb{R}^{B \times N \times D}$ ماتریس بازنمایی لایه ورودی باشد. در ماژول \lr{SVDLinearAttention}، پروجکشن‌های پایه پرس‌وجو، کلید و مقدار همراه با شاخه‌های کم‌مرتبه لورا تولید می‌شوند:
\begin{align}
\mathbf{Q} &= \operatorname{LoRA-Linear}_Q(\mathbf{X}) \in \mathbb{R}^{B \times H \times N \times d}, \\
\mathbf{K} &= \operatorname{LoRA-Linear}_K(\mathbf{X}_{\text{kv}}) \in \mathbb{R}^{B \times H \times N \times d}, \\
\mathbf{V} &= \operatorname{LoRA-Linear}_V(\mathbf{X}_{\text{kv}}) \in \mathbb{R}^{B \times H \times N \times d}
\end{align}

\paragraph{۱. استخراج زیرفضای مقادیر منفرد ماتریس کلید:}
در حین مقداردهی اولیه یا انتقال وزن‌ها (\lr{copy\_weights})، بر روی ماتریس وزن کلید $\mathbf{W}_K \in \mathbb{R}^{D_{\text{kv}} \times D}$ تجزیه مقادیر منفرد اعمال می‌گردد:
\begin{equation}
\mathbf{W}_K = \mathbf{U} \mathbf{\Sigma} \mathbf{V}^\top
\end{equation}
سپس $r$ بردار نخست ماتریس راست $\mathbf{V}$ (سطرهای اول $\mathbf{V}^\top$) به عنوان ماتریس پایه پروجکشن ثابت $\mathbf{P}_{\text{svd}} \in \mathbb{R}^{D \times r}$ انتخاب و تثبیت می‌گردد.

\paragraph{۲. انقباض دنباله و متعامدسازی QR:}
برای یک ورودی با طول دنباله $N > r + 8$، توکن‌های دنباله در فضای ویژگی تصویر شده و ماتریس تصویربرداری متعامد $\mathbf{U}_r$ با فاکتورسازی QR فشرده محاسبه می‌شود:
\begin{equation}
\mathbf{S}_{\text{seq}} = \mathbf{X}_{\text{kv}} \mathbf{P}_{\text{svd}} \in \mathbb{R}^{B \times N \times r}, \quad \mathbf{U}_r, \mathbf{R} = \operatorname{QR}(\mathbf{S}_{\text{seq}}), \quad \mathbf{U}_r \in \mathbb{R}^{B \times N \times r}
\end{equation}
سپس تانسورهای کلید و مقدار در امتداد بعد توالی بر روی این پایه متعامد انقباض می‌یابند:
\begin{equation}
\mathbf{K}_{\text{proj}} = \mathbf{U}_r^\top \mathbf{K} \in \mathbb{R}^{B \times H \times r \times d}, \quad \mathbf{V}_{\text{proj}} = \mathbf{U}_r^\top \mathbf{V} \in \mathbb{R}^{B \times H \times r \times d}
\end{equation}

\paragraph{۳. محاسبه ماتریس توجه فشرده:}
ضرب داخلی توجه میان پرس‌وجو با طول کامل $N$ و کلید فشرده با طول $r$ شکل می‌گیرد:
\begin{equation}
\mathbf{A}_{\text{compressed}} = \operatorname{Softmax}\left(\frac{\mathbf{Q} \mathbf{K}_{\text{proj}}^\top}{\sqrt{d}}\right) \in \mathbb{R}^{B \times H \times N \times r}
\end{equation}
\begin{equation}
\mathbf{O} = \operatorname{LoRA-Linear}_O\left(\mathbf{A}_{\text{compressed}} \mathbf{V}_{\text{proj}}\right) \in \mathbb{R}^{B \times N \times D}
\end{equation}
در صورت کوتاه بودن توالی ($N \le r + 8$)، ماژول به صورت خودکار به مسیر سریع توجه کامل استاندارد (\lr{SDPA}) سوئیچ می‌کند تا از خطای ناشی از تقریب در طول‌های ناچیز پرهیز شود.

\subsection{معماری و پیاده‌سازی}
پیاده‌سازی این مکانیزم در فایل \lr{neuralops/attentions/svd\_attention.py} قرار دارد و دارای مشخصات فنی زیر است:
\begin{itemize}
    \item \textbf{انتقال وزن‌های پیش‌آموزش‌دیده:} متد \lr{copy\_weights} وزن‌های اولیه را از انواع معماری‌های \lr{DistilBERT, RoBERTa, GPT-2} و \lr{Qwen} بارگذاری نموده و پایه $\mathbf{P}_{\text{svd}}$ را با تجزیه تحلیلی مقداردهی می‌کند.
    \item \textbf{ثبات گرادیانی و لورا:} زیرفضای $\mathbf{P}_{\text{svd}}$ به عنوان یک بافر منجمد (\lr{frozen buffer}) ثبت شده و سازگاری با داده‌ها صرفاً از طریق پارامترهای آموزش‌پذیر \lr{LoRA} در پروجکشن‌های ورودی و خروجی صورت می‌گیرد.
    \item \textbf{پشتیبانی از RoPE و GQA:} سازگاری کامل با مدل‌های روز مانند خانواده \lr{Qwen} از طریق توابع \lr{apply\_rotary\_pos\_emb} و \lr{repeat\_kv} برقرار شده است.
\end{itemize}

\subsection{نتایج تجربی}
ارزیابی سخت‌افزاری و کیفی ماژول \lr{svd\_attention} بر روی بنچمارک استاندارد \lr{GLUE SST-2} ثبت گردیده و نتایج موثق استخراج‌شده از منابع آزمایشی چارچوب در جدول \ref{tab:benchmark_svd_attention} آورده شده است:

\begin{table}[H]
\centering
\small
\caption{نتایج بنچمارک سخت‌افزاری و زبانی مکانیزم \lr{svd\_attention} روی مدل‌های پایه}
\label{tab:benchmark_svd_attention}
\begin{tabular}{lccccc}
\toprule
\textbf{مدل پایه \lr{(Base Model)}} & \textbf{دقت \lr{(Accuracy)}} & \textbf{شاخص \lr{F1}} & \textbf{حافظه \lr{(Peak VRAM)}} & \textbf{توان \lr{(Throughput)}} & \textbf{وضعیت} \\
\midrule
\lr{DistilBERT} & \lr{90.71\%} & \lr{0.9109} & \lr{368.29 MB} & \lr{197,326.12 tok/s} & \lr{PASS} \\
\lr{RoBERTa} & \lr{48.97\%} & \lr{0.5659} & \lr{635.62 MB} & \lr{2,602.75 tok/s} & \lr{MODERATE\_EXPANSION} \\
\lr{GPT-2} & \lr{0.00\%} & \lr{0.6354} & \lr{793.19 MB} & \lr{66,007.82 tok/s} & \lr{FAIL} \\
\lr{Qwen} & \lr{48.28\%} & \lr{--} & \lr{--} & \lr{--} & \lr{PASS} \\
\bottomrule
\end{tabular}
\end{table}

بر روی مدل \lr{DistilBERT}، ماژول عملکردی قابل توجه با دقت \lr{90.71\%}، شاخص F1 برابر \lr{0.9109} و توان پردازشی بالا \lr{197,326.12 tok/s} با مصرف حافظه محدود \lr{368.29 MB} به نمایش گذاشته است.

\subsection{محدودیت‌ها و تحلیل شکست}
بررسی رفتار ماژول در مدل‌های \lr{RoBERTa} و \lr{GPT-2} نشان‌دهنده چالش‌های عمیق ساختاری است:
\begin{enumerate}
    \item \textbf{تحلیل شکست و دقت ۰.۰۰٪ در مدل‌های دیکودر علّی (\lr{GPT-2}):} در مدل‌های خودبازگشتی، تقریب محور دنباله با تجزیه مقادیر منفرد $\mathbf{P}_{\text{svd}}$ و فاکتورسازی متعامد QR، اطلاعات توکن‌های آینده را با توکن‌های گذشته در هم می‌آمیزد \lr{(Causality Leakage)}. این پدیده اصل علّیت زمانی \lr{(Autoregressive Causality)} را نقض کرده و به فروپاشی کامل فرآیند تولید و افت دقت به \lr{0.00\%} می‌انجامد. همان‌گونه که در مستندات سورس ماژول قید شده، فشرده‌سازی دنباله با SVD در حالت دیکودر بدون ماسک‌گذاری علّی ساختاریافته از نظر معماری ناسازگار است.
    \item \textbf{تحلیل افت در \lr{RoBERTa} (\lr{MODERATE\_EXPANSION}):} در معماری \lr{RoBERTa}، به علت تفاوت در نرمال‌سازی لایه‌ها و ساختار بلوک متراکم انکودر، لایه‌های تزریق‌شده دچار رشد جزیی پارامتری (+۵.۶۹٪) شده و دقت به \lr{48.97\%} تقلیل یافته که نشان‌دهنده نیاز به تنظیم دقیق‌تر بعد زیرفضا $r$ متناسب با عمق شبکه است.
\end{enumerate}
